\section{Synthesis: information value is conditional, small, and policy-fragile}
\label{sec:synthesis}

\textbf{Three questions, asked with one instrument, returned the same answer.} Read
together, Q16, Q17, and Q18 rule out the monotone premise that organizes much of applied
market-microstructure learning. Reconstructing hidden structure did not improve the
decision signal (Q16). The best forecaster was not the best decider (Q17). And no fixed
observable configuration was universally sufficient (Q18). The common cause is visible
once the studies are aligned: value that is real at the level of a statistical target, a recovered queue size, a higher macro-F1, a deeper book, is small in magnitude and is
absorbed or reversed by the policy and cost layer that separates a prediction from a
decision. This is why we measure at the decision. A programme that had stopped at
reconstruction RMSE, at macro-F1, or at an averaged depth curve would have reported three
positive results. Measured at the decision, all three are negative.

\paragraph{The reusable object.} The contribution readers can depend on is the
\emph{matched information-intervention protocol}: freeze model, budget, dates, and policy. Vary only the observed information. Read the effect at the decision under an explicit cost
model. And refuse outcome-driven tweaks. The protocol is what turns ``does richer data
help?'' into a falsifiable measurement, and it transfers beyond these three questions to
any setting where a feed is bought on the promise of downstream value.

\paragraph{The released programme.} We release thirty curated experiment packages spanning
the three real-data studies, controlled simulator mechanisms, breakout/rebound cohorts,
temporal-persistence checks, and source-and-schema certification. Each package carries its
question, method, compact results, figures, provenance, and an explicit scientific-status
label, and each is laptop-reproducible from a license-clean smoke fixture. Raw market
dumps are never redistributed. The labels are load-bearing: a completed software check, a
valid scientific negative, an inconclusive run, and a blocked empirical objective are kept
distinct rather than rounded up to ``done.''

\paragraph{One honest open endpoint.} The strictest form of Q16, native-FIFO count
tightening on a fully admitted order-lifecycle cohort, is not yet computed. We located
the exact source (an open December 2025 Level-4 Hyperliquid archive of $195.3$ GB with
$744$ hourly files per order archive) and a bounded washout audit over $11$ complete hours
covered $32{,}633{,}820$ events, observing $11{,}328$ legacy order identifiers of which
$148$ were still first seen in hour $10$. Observed-lifecycle coverage reached
$99.92985\%$, but true live-book completeness remains unidentified, so we do not claim the
native-FIFO endpoint. We flag this rather than paper over it, and the decision-relevance
conclusion for Q16 stands on the reconstruction diagnostic, which needs no such admission.
